\section{Capitán en Jefe Lina María Marín Rodríguez --- Exdirectora Nacional de Bomberos de Colombia}

\subsection{Información básica}
La Capitán en Jefe Lina María Marín Rodríguez ejerció como Directora Nacional de Bomberos de Colombia desde junio de 2025. Su renuncia fue aceptada el 18 de agosto de 2026 y el nuevo director asumió el cargo el 7 de septiembre de 2026.

\subsection{Anotaciones}
\begin{itemize}
  \item Ser cordial, profesional y respetuoso.
  \item Agradecer su tiempo y explicar con claridad el propósito académico.
\end{itemize}

\subsection{Perfil profesional}
\begin{enumerate}
  \pregunta ¿Cuáles fueron sus principales responsabilidades durante su gestión como Directora Nacional de Bomberos de Colombia?
  \pregunta ¿Qué experiencia tuvo en la formulación, coordinación o seguimiento de estrategias frente a incendios?
  \pregunta ¿Cómo se articulaban las decisiones del nivel nacional con las necesidades de los cuerpos de bomberos territoriales?
  \pregunta ¿Cuáles fueron los retos más recurrentes que enfrentaron los cuerpos de bomberos para prevenir y atender incendios durante su gestión?
\end{enumerate}

\subsection{Planeación, información y escenarios}
\begin{enumerate}
  \pregunta ¿Qué herramientas de software, mapas de riesgo, modelos o sistemas de información se utilizaban para estimar el comportamiento de los incendios, cómo aportaban a las decisiones nacionales y cuáles eran sus principales limitaciones?
  \pregunta Desde su experiencia, ¿qué tipo de información habría sido más útil para fortalecer la toma de decisiones a nivel nacional frente a un incendio?
\end{enumerate}


% ---------------------------------------------------------
% CONTENIDO AÑADIDO
% ---------------------------------------------------------

\subsection{Coordinación y toma de decisiones a nivel nacional}

\begin{enumerate}

  \pregunta ¿Cómo llegaba la información desde los cuerpos de bomberos territoriales hasta el nivel nacional, qué tan homogénea era entre regiones y qué dificultades existían en su calidad, disponibilidad, actualización o estandarización?



  \pregunta Cuando ocurrían múltiples incendios simultáneamente, ¿cómo se obtenía una visión general de la situación nacional?


\end{enumerate}


\subsection{Anticipación del comportamiento de los incendios}

\begin{enumerate}
  \pregunta ¿Qué tan importante considera poder pasar de conocer dónde se encuentra actualmente un incendio a estimar dónde podría encontrarse algunas horas después y para qué decisiones a nivel nacional sería especialmente útil?


  \pregunta ¿Con qué anticipación debería obtenerse una estimación para que realmente pueda modificar una decisión?



  \pregunta ¿Qué utilidad tendría conocer no solo la trayectoria probable del incendio, sino también poblaciones, infraestructura o ecosistemas potencialmente afectados?

\end{enumerate}


\subsection{Datos e interoperabilidad}

\begin{enumerate}
  \pregunta ¿Qué instituciones generan actualmente datos relevantes para estudiar o anticipar incendios forestales y qué tan sencillo es obtenerlos y combinarlos?




  \pregunta ¿Considera necesario utilizar formatos o protocolos estandarizados para facilitar el intercambio de información entre instituciones?

\end{enumerate}


\subsection{Validación de una posible herramienta}

Antes de realizar las siguientes preguntas se explicará brevemente que el proyecto
estudia la posibilidad de generar previamente múltiples escenarios de propagación
de incendios bajo diferentes condiciones y utilizarlos posteriormente para
identificar con rapidez situaciones similares a una emergencia real.

\begin{enumerate}
  \pregunta Pensando en una herramienta de este tipo, ¿en qué momento de una emergencia y para qué decisiones concretas considera que podría aportar mayor valor?

  \pregunta ¿Qué información debería mostrar obligatoriamente para resultar útil?

  \pregunta ¿Qué información considera que podría resultar innecesaria o dificultar la interpretación?

  \pregunta ¿Qué factores podrían dificultar o facilitar que una institución adopte una herramienta de este tipo?
\end{enumerate}


\subsection{Viabilidad e implementación}

\begin{enumerate}
  \pregunta Si inicialmente se desarrollara únicamente un prototipo académico, ¿qué funcionalidad mínima debería tener para que pudiera ser evaluado por una entidad del sector?





  \pregunta Además de la atención de emergencias reales, ¿qué otros usos identifica para una herramienta basada en escenarios simulados, por ejemplo capacitación, prevención, investigación o planeación territorial?
\end{enumerate}


\subsection{Cierre}

\begin{enumerate}
  \pregunta Desde su experiencia, ¿cuál considera que es actualmente la mayor oportunidad para mejorar tecnológicamente la gestión de incendios forestales en Colombia y, si el proyecto solo pudiera resolver correctamente un problema, cuál recomendaría priorizar?

  \pregunta ¿Qué aspecto de una solución como la planteada considera que deberíamos investigar con mayor profundidad antes de desarrollarla?

  \pregunta ¿Qué otras personas, entidades o perfiles profesionales considera que deberíamos entrevistar para comprender mejor este problema?

  \pregunta ¿Conoce algún proyecto, sistema, plataforma o iniciativa nacional o internacional que considere importante revisar como referencia para nuestro proyecto?
\end{enumerate}
